\documentclass[12pt]{article}
\usepackage[utf8]{inputenc}
\usepackage[T1]{fontenc}
\usepackage{hyperref}
\title{Sawila Language Resources for K-12 Curriculum Development}
\author{Compiled by Claude Sonnet 5.5 (claude-sonnet-5-5) for Andr\'as Kornai}
\date{}
\begin{document}
\maketitle

\section{Introduction}
Sawila (ISO 639-3 swt; Glottocode sawi1256), also called Tanglapui, is a Papuan language of Alor, Indonesia, in the Timor-Alor-Pantar family. Wikipedia gives about 3,000 speakers (2014). Dialects named there are Sawila proper, Lona, Salimana, Lalamana and Sileba. Speakers call the language Manata. Sawila is grouped with Kula as the Tanglapui languages. The only substantial resource found is a scholarly grammar sketch.

\section{Primary Educational Resources}
No primers, textbooks or school programmes for Sawila were found in search results.

\section{Linguistic and Reference Materials}
\begin{itemize}
\item Kratochv\'il, Fran\v{t}i\v{s}ek (2014), ``Sawila'', in Schapper (ed.), \emph{The Papuan Languages of Timor, Alor and Pantar}, vol.\ 1, De Gruyter, pp.\ 351--438 (sketch grammar; DOI 10.1515/9781614515241.351). Bibliographic record: \url{https://glottolog.org/resource/reference/id/319022}
\item Wikipedia, ``Sawila language'': \url{https://en.wikipedia.org/wiki/Sawila_language}
\item Wikipedia, ``Tanglapui languages'': \url{https://en.wikipedia.org/wiki/Tanglapui_languages}
\item A reference listing mentions a Sawila dictionary by Kratochv\'il, Bantara and Malaikosa. I did not see the dictionary itself, and the listing is in this Leiden repository item: \url{https://scholarlypublications.universiteitleiden.nl/access/item%3A3719511/view}
\end{itemize}

\section{Religious and Community Resources}
None found in search results.

\section{Digital and Audio Resources}
None found in search results. The University of Zurich New Guinea World page on the Alor-Tip languages gives family context only: \url{https://newguineaworld.linguistik.uzh.ch/families/trans-new-guinea/berau-gulf/west-bomberai/timor-alor-pantar/alor-pantar/east-alor/alor-tip/start}

\section{Curriculum Development Context}
The language is not known to be taught in schools. The Kratochv\'il grammar is the main basis for any orthography or teaching grammar. Educational materials would need to be created with the community.

\end{document}